% ASSERT_CONVENTION: natural_units=dimensionless, metric_signature=mostly_minus,
%   jordan_product=(1/2)(ab+ba), octonion_basis=fano_e1e2=e4,
%   complex_structure=u_equals_e7, potential=minus_log_det,
%   cone_metric=hess(minus_log_det), cubic_norm_det=bulk_geometry_verification_det_3,
%   curvature_engine=totaro_closed_form, arithmetic=exact_over_Q
%
% Phase 70.1 (A0' -- REVISE A0), Plan 01: select the physical spacetime metric --
%   the Riemannian cone-Hessian g_X = Hess(-log det_3) vs the construction-(ii)
%   bridge g = eta + h -- and restate the route thesis accordingly, for the v17.0
%   milestone "Gravity as Intrinsic Curvature of the h_3(O) Bulk Geometry".
%
% This phase does NOT rebuild the engine, does NOT recompute the Totaro Riemann, and
% does NOT run a numerical sweep. The engine, det_3 (SSOT), the Totaro curvature
% method, and the {0,-1,-1,-1} Ricci spectrum are ALREADY certified exact-over-Q
% (code/bulk_geometry_verification.py; reproduced for THIS phase by
% .gpd/phases/70.1-.../70.1-01-metric-selection.py, final line
% METRIC_SELECTION_INPUTS_OK, exit 0). We reuse, we do not re-derive.
%
% STATUS: Tasks 1-3 written (headline + Einstein verdict + the two interpretation
%   analyses). Task 4 is a BLOCKING human-decision checkpoint (ratify the thesis
%   restatement: cone-Hessian / eta+h / escalate). Tasks 5 (selection verdict) and 6
%   (consequence propagation to Phases 72/73) are written ONLY AFTER ratification.
%
% References (project-consistent tags; cf. 72-matter-sourcing.tex, 70-signature-bridge.tex):
%   [TOT04] Totaro, "The curvature of a Hessian metric," arXiv:math/0401381
%           (Int. J. Math. 15 (2004) 369-391), Cor. 2.3 closed form; constant
%           sectional curvature -d^2/4 for the rank-1 cone slice. PLAN-LOCAL id
%           ref-totaro (the curvature reference lives in CONVENTIONS.md, not in the
%           project contract reference list). ALREADY in the warm engine -- cite the
%           method, do NOT recompute.
%   [FK94]  Faraut & Koranyi, Analysis on Symmetric Cones (OUP 1994), Ch. II-IV:
%           g_X = Hess(-log det) positive-definite / Riemannian (WHY a Riemannian ->
%           Lorentzian bridge is forced at all).
%   [Besse] A. L. Besse, Einstein Manifolds (Springer 1987): Einstein condition
%           Ric = lambda g, lambda = R/n (n>=3 => lambda const); traceless-Ricci
%           S = Ric - (R/n)g, Einstein <=> S == 0; metric cones (Sec 9 / Ch. 1).
%   [52KKT] derivations/52-kkt-spacetime.tex + 52-observer-uniqueness.tex:
%           h_2(C_u) ~ R^{3,1}, det_2 = x0^2-x1^2-x2^2-x3^2, signature (1,3)
%           mostly-minus, Minkowski coords x0=(beta+gamma)/2, x3=(beta-gamma)/2,
%           timelike x_0 ~ (beta+gamma) = (1,1,0,0); H^3 = SL(2,C)/SU(2).
%   [h3o]   ~/repos/blog/research/qualia-fixed-point/h3o_tower.py: the corrected
%           cubic-norm cross-term order (the warm engine det_3 is byte-identical to
%           the certified SSOT; the buggy (x1 x2)x3 / octonion_algebra.py order is
%           banned -- fp-wrong-cross-term).
%   [70-02] derivations/70-signature-bridge.tex + 70-02-SUMMARY.md: construction-(ii)
%           bridge g = eta + h with h := Hess - Hess|center; Minkowski-reduction gate
%           flagged tautological-by-construction (Note B); construction (i) REJECTED.
%   [72-01] derivations/72-matter-sourcing.tex (Sec obstruction, Sec adjudication,
%           Sec phase70-task): the VALD-04 finding (the non-Einstein M=0 vacuum IS
%           the result, not a reference choice); reconciliation (a) REJECTED as
%           fp-relabel, (b) CONDITIONAL; the two-outcome framing this phase resolves.
%   [FP-cb] Fierz-Pauli / linearized-Einstein on a curved background: the massless
%           spin-2 gauge (vector) symmetry extends only to an Einstein background
%           (Ric proportional to g). arXiv:2503.11304 (EPJC 2025); arXiv:1109.1382.
%           Background-/dimension-agnostic; see Task 3.

\section{Phase 70.1-01: selecting the physical spacetime metric (cone-Hessian vs $\eta+h$)}

\subsection{Conventions and scope (carried verbatim from the v17.0 lock)}

All decisive quantities are EXACT over $\mathbb{Q}$ (no float on any verdict; ranks and
eigenvalues via \texttt{sympy.Matrix}, \texttt{.eigenvals()}). The cone metric is the
Hessian of the potential $\Phi=-\log\det_3$, with $\det_3$ the certified Freudenthal
cubic norm (engine SSOT, cross-term order $2\,\mathrm{Re}((x_2 x_1)x_3)$;
\texttt{octonion\_algebra.py} banned -- \texttt{fp-wrong-cross-term}). The Riemann
tensor is the Totaro [TOT04] closed form
\begin{equation}
  R_{ijkl} = -\tfrac14\, g^{pq}\big(C_{jlp}C_{ikq}-C_{ilp}C_{jkq}\big),
  \qquad C_{ijk}=\Phi_{,ijk},
  \label{eq701:totaro}
\end{equation}
with $\mathrm{Ric}_{jl}=g^{ik}R_{ijkl}$, $R=g^{jl}\mathrm{Ric}_{jl}$ (already implemented
and certified in the warm engine -- not recomputed here). The spacetime sub-slice is
$h_2(C_u)$, engine-native indices $\{1,2,3,10\}=(\beta,\gamma,p,q)$, mostly-minus
$(-,+,+,+)$ on the slice, Riemannian positive-definite in the bulk; the center is
$X_{\mathrm{bg}}=I/3$. The Riemann/Ricci sign is pinned NEGATIVE and constant
(cone-Hessian-slice $K(p,q)=-\tfrac12$, round $H^3$ $K=-1$; the factor-of-$2$ is BENIGN
-- the load-bearing fact is the negative SIGN; we do NOT force $-1$:
\texttt{fp-assume-einstein}).

\paragraph{The one load-bearing question.} Phase 72's first-result gate computed the
$M=0$ cone-Hessian vacuum exactly over $\mathbb{Q}$ and found it is the non-Einstein
static product $\mathbb{R}_{\text{time}}\times H^3$ ([72-01], Sec.~obstruction). The
Riemannian cone-Hessian $g_X=\mathrm{Hess}(-\log\det_3)$ (the route thesis) and the
Lorentzian construction-(ii) bridge $g_{\mu\nu}=\eta_{\mu\nu}+h_{\mu\nu}$ ([70-02])
therefore \emph{disagree} about the $M=0$ geometry (cone-Hessian $\to
\mathbb{R}\times H^3$; $\eta+h\to$ flat Minkowski), and that disagreement IS the
underdetermination of the A0 signature bridge. This plan selects which object is the
physical spacetime metric and restates the thesis at its TRUE strength -- OR, if neither
restatement is defensible, escalates (binding backtracking trigger). The decisive inputs
below are reproduced by \texttt{70.1-01-metric-selection.py} (imports the warm engine;
\texttt{main()} not run; final line \texttt{METRIC\_SELECTION\_INPUTS\_OK}; exit $0$,
$\sim$3.5\,s).

\subsection{Headline: the $M=0$ cone-Hessian vacuum is the non-Einstein static
            product $\mathbb{R}_{\text{time}}\times H^3$ (reproduced, not re-derived)}
\label{sec701:headline}

Reproduced EXACT over $\mathbb{Q}$ from the warm engine ([72-01] eq:ricci-center;
\texttt{70.1-01-metric-selection.py} steps 2--5):
\begin{equation}
  g = \mathrm{diag}(9,9,18,18)\ \ (\det g = 26244),\qquad
  \mathrm{Ric} =
  \begin{pmatrix} -\tfrac92 & \tfrac92 & 0 & 0\\
                  \tfrac92 & -\tfrac92 & 0 & 0\\
                  0 & 0 & -18 & 0\\
                  0 & 0 & 0 & -18\end{pmatrix},
  \label{eq701:ricci-center}
\end{equation}
\begin{equation}
  g^{-1}\mathrm{Ric}\ \text{has eigenvalues}\ \{\,0,-1,-1,-1\,\}\
  (\text{mult } 0{:}1,\ {-1}{:}3),\qquad R = -3.
  \label{eq701:spectrum}
\end{equation}
The coordinate sectional curvatures name the structure:
\begin{equation}
  K(e_0,e_1) = 0\ \ (\text{timelike/dilation 2-plane FLAT}),\qquad
  K(p,q) = K(e_2,e_3) = -\tfrac12\ \ (\text{$H^3$ base}),
  \label{eq701:sectional}
\end{equation}
the $H^3$ value identical to the engine's certified
\texttt{h3\_cone\_hessian\_benchmark} ($K_{\text{value}}=-\tfrac12$; round-$H^3$
cross-check $K_{\text{round}}=-1=2\cdot(-\tfrac12)$, the benign factor-of-$2$). The
eigenvalue-$0$ eigenvector is $(1,1,0,0)=\beta+\gamma$; under the Lorentzian reading
$h_2(C_u)\simeq\mathbb{R}^{3,1}$ with $\det_2$ Minkowski form
$\beta\gamma/3-p^2/3-q^2/3$ ([52KKT]), this is the \textbf{timelike $x_0$} direction (the
$\beta\gamma/3$ block is the hyperbolic $1{+}1$ part; $(1,1)$ its timelike eigenvector).
So the flat factor of the cone is \emph{time}, and the $M=0$ vacuum is the static product
$\mathbb{R}_{\text{time}}\times H^3_{\text{space}}$. This spectrum is the HEADLINE: it is
the VALD-04 finding (negative-result-is-success), reported as such and never papered over.

\paragraph{Inherited and confirmed (carry-forward gates).} The construction-(ii)
Minkowski reduction (test-minkowski-reduction) is satisfied by the warm engine:
$g(\text{center},M{=}0)=\eta$ exactly, since $h:=\mathrm{Hess}-\mathrm{Hess}|_{\text{center}}=0$
at the center by construction ([70-02] Note B; \texttt{70.1-01-metric-selection.py}
step 6: \texttt{offcenter\_slice\_metric}$\{\}$ gives $h=\mathbf{0}_{4\times4}$,
$g=\eta_{\mathrm{bg}}$; \texttt{minkowski\_reduction()} residual $=0$, Sylvester minors
$[1,-1,1,-1]\Rightarrow$ signature $(1,3)$). The corrected cubic-norm cross-term
(test-cross-term-association) is confirmed via the engine SSOT: $\det_3$ minus the
block-diagonal norm $\alpha\beta\gamma-\alpha|x_1|^2-\beta|x_2|^2-\gamma|x_3|^2$ equals
\emph{exactly} $2\,\mathrm{Re}((x_2 x_1)x_3)$, with $\alpha,\beta,\gamma$ absent from the
triple (\texttt{70.1-01-metric-selection.py} step 8); the buggy $(x_1 x_2)x_3$ /
\texttt{octonion\_algebra.py} order is excluded ([h3o]; \texttt{fp-wrong-cross-term}).

\subsection{Einstein-condition verdict: the cone-Hessian $M=0$ vacuum is NOT a GR
            $\Lambda$-vacuum}
\label{sec701:einstein}

An Einstein space satisfies $\mathrm{Ric}_{\mu\nu}=\lambda\,g_{\mu\nu}$ with
$\lambda=R/n$ (for $n\ge3$, $\lambda$ is necessarily constant) [Besse]. Equivalently the
traceless Ricci tensor
\begin{equation}
  S_{\mu\nu} := \mathrm{Ric}_{\mu\nu}-\tfrac{R}{n}\,g_{\mu\nu}
  \qquad\text{satisfies}\qquad
  \text{Einstein}\iff S_{\mu\nu}\equiv 0
  \label{eq701:tracelessricci}
\end{equation}
[Besse; the Einstein criterion]. For $n=4$ and $R=-3$ the Einstein value is
\begin{equation}
  \lambda = \frac{R}{4} = -\frac34,
  \label{eq701:einstein-value}
\end{equation}
so an Einstein space in $n=4$ with $R=-3$ would require \emph{all} Ricci-endomorphism
eigenvalues to equal $-\tfrac34$. The computed spectrum~\eqref{eq701:spectrum} is
$\{0,-1,-1,-1\}$, which contains \textbf{no} eigenvalue equal to $-\tfrac34$. Therefore
the cone-Hessian $M=0$ vacuum is NOT Einstein, hence NOT a GR $\Lambda$-vacuum. The
direct exact-over-$\mathbb{Q}$ witness is the nonzero traceless Ricci tensor
(\texttt{70.1-01-metric-selection.py} step 4):
\begin{equation}
  S = \mathrm{Ric}-\tfrac{R}{4}\,g
    = \mathrm{Ric}+\tfrac34\,g
    =
  \begin{pmatrix} \tfrac94 & \tfrac92 & 0 & 0\\
                  \tfrac92 & \tfrac94 & 0 & 0\\
                  0 & 0 & -\tfrac92 & 0\\
                  0 & 0 & 0 & -\tfrac92\end{pmatrix}
  \;\neq\; \mathbf{0},
  \qquad \operatorname{tr}_g S = g^{\mu\nu}S_{\mu\nu} = 0
  \label{eq701:S}
\end{equation}
(genuinely traceless yet nonzero -- the exact signature of a non-Einstein space). The
$n=4$ Ricci decomposition reconstruction residual $R_{ijkl}-(\mathrm{Scal}+E+\mathrm{Weyl})=0$
holds exactly over $\mathbb{Q}$, confirming $S\neq0$ is a genuine curvature fact, not a
decomposition artifact.

\paragraph{Signature-independence of the verdict.} The eigenvalues
$\{0,-1,-1,-1\}$ are invariants of the \emph{$(1,1)$} Ricci endomorphism
$g^{-1}\mathrm{Ric}$, not of the raw $(0,2)$ Ricci components. A change of frame
(congruence) $P$ acts on the $(0,2)$ tensors by $g\mapsto P^{\!\top}gP$,
$\mathrm{Ric}\mapsto P^{\!\top}\mathrm{Ric}\,P$, so the endomorphism transforms by
\emph{similarity},
\begin{equation}
  g^{-1}\mathrm{Ric}\;\longmapsto\;(P^{\!\top}gP)^{-1}(P^{\!\top}\mathrm{Ric}\,P)
  = P^{-1}\,(g^{-1}\mathrm{Ric})\,P,
  \label{eq701:similarity}
\end{equation}
and its eigenvalues are invariant. Converting between the $(-,+,+,+)$ and $(+,-,-,-)$
labelings is such a congruence, so it does NOT change the verdict; this is verified
exact over $\mathbb{Q}$ (\texttt{70.1-01-metric-selection.py} step 7: a generic
invertible rational $P$ leaves the spectrum $\{0,-1,-1,-1\}$ fixed).

\emph{Caution (why the invariant is the $(1,1)$ endomorphism, index raised).} A naive
\emph{overall} metric sign flip $g\mapsto-g$ is \textbf{not} a congruence and is NOT the
same as a signature change: it leaves the $(0,2)$ tensor $\mathrm{Ric}_{\mu\nu}$
unchanged but negates the $(1,1)$ endomorphism $g^{-1}\mathrm{Ric}$ (eigenvalues
$\{0,-1,-1,-1\}\mapsto\{0,+1,+1,+1\}$, verified over $\mathbb{Q}$ in step 7). This is
exactly the contract's disconfirming observation -- ``a sign-dependent verdict means the
$(0,2)$ Ricci components were used without raising the index.'' The verdict we report is
the \emph{similarity-invariant} $(1,1)$ spectrum~\eqref{eq701:spectrum}, computed with
the index raised; no verdict in this phase uses the raw $(0,2)$ entries or an overall
scalar flip.

\paragraph{Explicit rejection of the ``which $\Lambda$ reference?'' baseline.} No
baseline subtraction turns a non-Einstein vacuum into an Einstein one. Asking ``which
$\Lambda$ reference?'' presupposes a maximally-symmetric vacuum that this construction
does not produce ([72-01], Sec.~adjudication; \texttt{fp-assume-einstein}). The spectrum
$\{0,-1,-1,-1\}$ IS the result -- headline it, do not paper it over. We insert no factor
to force $\mathrm{Ric}\propto g$ or to force flatness, and we keep $K(p,q)=-\tfrac12$
(no factor of $2$ to ``match'' $-1$).

\paragraph{Why a metric cone is never Einstein (corroborative, not load-bearing).} The
dim-$4$ cone-Hessian sub-slice is a metric cone $C(H^3)=\mathbb{R}_{>0}\times H^3$: a
flat radial (dilation) line over the constant-curvature $H^3$ base. A non-trivial metric
cone in dimension $n$ has a radial $0$-eigenvalue plus shifted base eigenvalues, here
$\{0,-1,-1,-1\}$, and is never Einstein [Besse Sec.~9; Gallot]. This EXPLAINS why the
spectrum has the form it does; the exact-over-$\mathbb{Q}$
computation~\eqref{eq701:spectrum}--\eqref{eq701:S} \emph{is} the proof. (The radial
$0$-direction is at once the Riemannian scale modulus $R(X)=R(2X)$ found in Phase 71 and
-- under the Lorentzian reading -- the timelike $x_0$; the cone's flat factor is time.)

\subsection{Interpretation analysis I: linearized spin-2 on a fixed non-Einstein
            background -- the option-(1) defensibility test}
\label{sec701:spin2}

Suppose the cone-Hessian is named the physical metric (option (1)). Its $M=0$ vacuum is
the provably non-Einstein $\mathbb{R}\times H^3$~\eqref{eq701:spectrum}. What, then, can
a perturbation $h_{\mu\nu}$ on this background legitimately be called?

\paragraph{Standard result (background-/dimension-agnostic).} The massless Fierz--Pauli
/ linearized-Einstein (Lichnerowicz) operator propagates a \emph{consistent massless
spin-2 field} only when its linearized gauge (vector) symmetry
$\delta h_{\mu\nu}=\nabla_\mu\xi_\nu+\nabla_\nu\xi_\mu$ is preserved; that gauge symmetry
of the massless equations extends to a curved background \emph{if and only if the
background is Einstein} ($\mathrm{Ric}_{\mu\nu}\propto g_{\mu\nu}$, including the
Ricci-flat and (A)dS cases) [FP-cb]. On a generic non-Einstein background the would-be
gauge variation fails to leave the linearized operator invariant (the obstruction is
proportional to the traceless Ricci $S_{\mu\nu}\neq0$), so there is no consistent
massless spin-2 gauge symmetry. This theorem is \emph{background- and
dimension-agnostic}: it is not a de Sitter/AdS-specific statement, and it applies to any
non-Einstein Lorentzian background, a metric cone $\mathbb{R}\times H^3$ included.

\paragraph{Confidence and the cone-specific caveat.} This is a MEDIUM-confidence
\emph{application} of a HIGH-confidence theorem. The theorem (gauge invariance
$\Rightarrow$ Einstein background) is multiply confirmed and background-agnostic; but no
reference treats massless spin-2 on a \emph{metric cone} $\mathbb{R}\times H^3$
specifically. We therefore rely on the theorem's stated generality, not on a
cone-specific citation, and flag this as the weakest anchor of the phase.

\paragraph{Strongest honest option-(1) claim (NAMED at true strength).} On the
non-Einstein $\mathbb{R}\times H^3$ the strongest defensible statement is a
\textbf{fixed-background linear matter-response}: a wave operator / Green's function for
$h_{\mu\nu}$ sourced by matter on a \emph{fixed} (non-dynamical, non-Einstein)
background, carrying \emph{no} linearized gauge symmetry and \emph{no} self-consistency
(no back-reaction closing into a gravitational sector). It must be named exactly that. It
is NOT ``GR derived,'' NOT ``a graviton,'' and NOT ``a self-consistent gravitational
spin-2 sector'' (\texttt{fp-overclaim-spin2}). The structure \emph{looks} like linearized
gravity (curved background $+$ perturbation), and ``gravity'' is the hoped-for answer --
which is precisely why the over-claim must be refused and the weaker object named at its
true strength.

\paragraph{Statable-vs-escalate recommendation (tested against the binding trigger).}
The EXPECTED and recommended conclusion is: a fixed-background matter-response \emph{is}
statable (as named above); it is \emph{not} a gravitational spin-2 sector. A
strictly-weaker-but-honest claim is a SUCCESS, not an escalation trigger. The escalation
trigger fires only on a \emph{conjunction} (see Sec.~\ref{sec701:checkpoint}): option (1)
fails \emph{and} option (2) fails. The open risk attached to option (1) alone is Open
Question 1: \emph{is a fixed-background wave operator without gauge invariance a
physically meaningful ``matter-response'' object at all, or is the absence of gauge
invariance fatal even to that reading?} If the obstruction is judged fatal even to the
matter-response object, option (1) fails -- and then, only if option (2) also fails, does
the route's gravitational interpretation become NOT established and the trigger fire.

\subsection{Interpretation analysis II: the $\eta+h$ circularity -- the option-(2)
            tripwire}
\label{sec701:circularity}

Suppose instead the construction-(ii) bridge $g_{\mu\nu}=\eta_{\mu\nu}+h_{\mu\nu}$ is
named the physical metric (option (2)). By construction
\begin{equation}
  h_{\mu\nu} := \big[\mathrm{Hess}(-\log\det_3)\ \text{restricted to } V_0\big]
              - \big[\text{its value at } (M{=}0,\ \text{center } I/3)\big],
  \label{eq701:hdef}
\end{equation}
so $h_{\mu\nu}\equiv0$ at the center identically ([70-02] Note B; reproduced in
\texttt{70.1-01-metric-selection.py} step 6 as $h=\mathbf{0}_{4\times4}$). Hence
$g(\text{center})=\eta$ is flat \emph{tautologically}, and $\Lambda=0$ is INSERTED by the
center-subtraction, \emph{not derived}. The Minkowski-reduction gate
(test-minkowski-reduction) is therefore tautological-by-construction: it confirms the
construction is implemented correctly and the signature is $(1,3)$, but it is NOT
independent evidence of an uncontaminated, dynamically-flat vacuum.

\paragraph{Exact tripwire string carried to Phase 73.} The following statement is carried
verbatim into Phase 73's per-equation circularity audit:
\begin{quote}\itshape
``$\Lambda=0$ at the center is inserted by the centered subtraction
$h:=\mathrm{Hess}-\mathrm{Hess}|_{\text{center}}$ ($h\equiv0$ at center by construction),
NOT derived; Phase 73 must audit whether any downstream Einstein-structure `fit' silently
re-uses this inserted flatness.''
\end{quote}

\paragraph{Both analyses feed the selection and the checkpoint.} Option (1) is defensible
iff the fixed-background matter-response survives as a meaningful object
(Sec.~\ref{sec701:spin2}, Open Question 1); option (2) is defensible iff the inserted
$\Lambda=0$ circularity can be honestly \emph{carried} as a tripwire
(Sec.~\ref{sec701:circularity}) rather than hidden. Both verdicts are inputs to the
human-decision checkpoint (Sec.~\ref{sec701:checkpoint}); neither is resolved here.

\subsection{Decision checkpoint: ratify the thesis restatement (PENDING)}
\label{sec701:checkpoint}

This is the BLOCKING human-decision point. The thesis restatement redefines the
milestone's central claim and re-scopes Phases 72/73, so under balanced autonomy the
adopted outcome must be ratified before the consequence-propagation is written. The
decision evidence assembled above is:
\begin{itemize}
  \item the reproduced exact-over-$\mathbb{Q}$ spectrum $\{0,-1,-1,-1\}$, $R=-3$,
        $S\neq0$~\eqref{eq701:spectrum}--\eqref{eq701:S}; $K(e_0,e_1)=0$,
        $K(p,q)=-\tfrac12$~\eqref{eq701:sectional} ($\mathbb{R}\times H^3$); $\eta+h$
        $h\equiv0$-at-center~\eqref{eq701:hdef}; signature-independence~\eqref{eq701:similarity};
  \item the Einstein verdict (Sec.~\ref{sec701:einstein}): NOT a GR $\Lambda$-vacuum;
  \item option (1) strongest honest claim $=$ fixed-background linear matter-response,
        NOT self-consistent spin-2 (Sec.~\ref{sec701:spin2}); option (2) $\eta+h$
        flatness is $\Lambda=0$-inserted, a circularity tripwire (Sec.~\ref{sec701:circularity});
  \item the binding backtracking trigger: escalate (do NOT force a choice) IFF option (1)
        does not admit even a fixed-background matter-response reading AND option (2)
        cannot be justified without circularity.
\end{itemize}
The admissible outcomes are exactly three: \textbf{(1)} cone-Hessian (thesis $=$
linearized matter-response on the fixed non-Einstein $\mathbb{R}\times H^3$);
\textbf{(2)} $\eta+h$ bridge (thesis $=$ bridge-metric-is-gravity, $\Lambda=0$-inserted
tripwire to Phase 73); \textbf{(escalate)} gravitational interpretation NOT established.
Sections~\ref{sec701:selection} (selection verdict) and the consequence propagation are
written only after the ratified outcome is recorded here.

% --- TASK 5 (selection verdict) and TASK 6 (consequence propagation) are appended
% --- below ONLY after the Task-4 human ratification of {1, 2, escalate}.
% \subsection{Selection verdict and restated thesis}
% \label{sec701:selection}
